\documentclass[11pt]{article}

\usepackage[a4paper,margin=1in]{geometry}
\usepackage[T1]{fontenc}
\usepackage[utf8]{inputenc}
\usepackage{lmodern}
\usepackage{amsmath,amssymb}
\usepackage{booktabs}
\usepackage{array}
\usepackage{float}
\usepackage{graphicx}
\usepackage{hyperref}
\usepackage{enumitem}

\title{Continuous-Action Warehouse Redistribution with PPO\\[4pt]
\large Technical Report — Stable Baselines 3 Implementation}
\author{}
\date{}

\begin{document}
\maketitle

%% ======================================================================
\section{Problem Setting}
%% ======================================================================

We consider a multi-warehouse inventory redistribution problem formulated as a Markov Decision Process (MDP). At each discrete time step $t$, an agent observes the inventory levels and stochastic demand at $N$ warehouses, then decides how much stock to ship between every pair. The objective is to minimise the total operating cost over an episode of $H=50$ steps.

\subsection{State, Action, and Observation}

The environment state is $s_t = (\mathbf{I}_t, \mathbf{D}_t)$, where $\mathbf{I}_t \in \mathbb{R}^N_{\geq 0}$ is the inventory vector and $\mathbf{D}_t \in \mathbb{R}^N_{\geq 0}$ is the demand vector. Demand at each warehouse is drawn independently from a Poisson distribution with mean $\mu = 10$.

The action is a transport matrix $T \in \mathbb{R}^{N \times N}_{\geq 0}$, where $T_{ij}$ denotes units shipped from warehouse $i$ to warehouse $j$. The environment enforces feasibility constraints via a projection operator:
\begin{equation}
    T_{ii} = 0, \qquad \sum_j T_{ij} \leq I_i \quad \forall\, i.
\end{equation}

The observation vector fed to the policy network is a normalised concatenation with a temporal feature:
\[
    o_t = \left[\frac{\mathbf{I}_t}{I_{\max}},\;\frac{\mathbf{D}_t}{3\mu},\;\frac{t}{H}\right] \in \mathbb{R}^{2N+1}.
\]
Inventory is divided by $I_{\max}=500$, demand by $3\mu=30$ (to keep values in $[0,1]$), and the time feature $t/H$ lets the policy distinguish early-episode (can afford to wait) from late-episode (must act now) situations.

\subsection{Reward Design}

The reward at each step is a weighted sum of four normalised penalties:
\begin{equation}\label{eq:reward}
    r_t = -\frac{C_{\text{transport}}}{\hat{C}_{\text{transport}}}
          - \lambda_d \,\frac{\sum_i U_i}{\hat{U}}
          - \lambda_r \,\frac{C_{\text{replenish}}}{\hat{C}_{\text{replenish}}}
          - \frac{C_{\text{overflow}}}{\hat{U}}
\end{equation}
where:
\begin{itemize}[nosep]
    \item $C_{\text{transport}} = \sum_{ij} c_{ij}\, T_{ij}$ is the total shipping cost,
    \item $U_i = \max(0,\, D_i - I_i^{\text{post}})$ is unmet demand at warehouse $i$,
    \item $C_{\text{replenish}} = c_{\text{ext}}\sum_i U_i$ is the cost if an external supplier fills shortages at $c_{\text{ext}}=50$/unit,
    \item $C_{\text{overflow}} = \sum_i \max(0,\, I_i - I_{\text{cap}})$ penalises inventory exceeding the per-warehouse cap $I_{\text{cap}} = 10\mu$,
    \item $\hat{C}_{\text{transport}},\, \hat{U},\, \hat{C}_{\text{replenish}}$ are worst-case normalisation constants derived from the environment parameters (not hand-tuned).
\end{itemize}

Each normalisation constant is computed from the actual cost matrix and demand configuration:
\begin{align*}
    \hat{C}_{\text{transport}} &= c_{\max} \cdot N(N-1) \cdot I_{\max}, \\
    \hat{U} &= N \cdot \mu, \\
    \hat{C}_{\text{replenish}} &= c_{\text{ext}} \cdot N \cdot \mu.
\end{align*}
Here $c_{\max} = \max_{ij} c_{ij}$ is derived from the actual cost matrix. This ensures that each reward term lies in approximately $[0, 1]$ before weighting, so the total per-step reward $r_t \in [-4, 0]$ regardless of scenario size. Without this normalisation, the raw reward ranged from $-900$ to $0$ per step, which destabilised value function learning.

The weighting hyperparameters $\lambda_d = 20$ and $\lambda_r = 1$ make unmet demand the dominant signal, encouraging the policy to prioritise service level over transport cost minimisation.

\subsection{Cost Matrix Structure}

The cost matrix follows a hub-and-spoke topology with $\lceil 0.2 N \rceil$ hub warehouses:
\begin{center}
\begin{tabular}{lc}
\toprule
Route type & Cost range \\
\midrule
Hub $\leftrightarrow$ Hub & $\mathcal{U}(0.1, 0.3)$ \\
Hub $\leftrightarrow$ Spoke & $\mathcal{U}(0.3, 0.8)$ \\
Spoke $\leftrightarrow$ Spoke & $\mathcal{U}(1.0, 2.0)$ \\
\bottomrule
\end{tabular}
\end{center}
For $N=4$ (small scenario), this yields 1 hub and 3 spokes. The cost structure creates a natural incentive to route stock through hubs, which a successful policy should discover.

\subsection{Evaluation Protocol}

Every 50 training episodes, the current policy is evaluated deterministically (no exploration noise) over 10 episodes in a separate environment instance. We report:
\begin{itemize}[nosep]
    \item \textbf{Episode reward}: cumulative $\sum_{t=0}^{H-1} r_t$,
    \item \textbf{Transport cost}: cumulative $\sum_t C_{\text{transport}}^{(t)}$,
    \item \textbf{Demand satisfaction}: mean $\frac{\sum_i \min(I_i^{\text{post}}, D_i)}{\sum_i D_i}$ over all steps,
    \item \textbf{Inventory std}: standard deviation of time-averaged inventory across warehouses (measures imbalance).
\end{itemize}

%% ======================================================================
\section{Methodology: PPO with Stable Baselines 3}
%% ======================================================================

\subsection{Why PPO}

Proximal Policy Optimisation (PPO)~\cite{schulman2017ppo} is a policy gradient method that constrains each update step to prevent catastrophic policy changes. For a continuous-action problem with an $N^2$-dimensional action space (16 dimensions for $N=4$), PPO is a natural choice: it handles continuous actions natively via a Gaussian policy, is on-policy (which avoids the deadly triad of off-policy + function approximation + bootstrapping), and is the most widely validated algorithm in Stable Baselines 3.

We use \textbf{Stable Baselines 3} (SB3), a PyTorch-based library of vetted RL implementations. SB3 provides a production-quality PPO with Generalised Advantage Estimation (GAE), vectorised environments, and reward normalisation — components that would take significant effort to implement correctly from scratch.

\subsection{PPO Algorithm}

PPO maximises the clipped surrogate objective:
\begin{equation}\label{eq:ppo}
    L^{\text{CLIP}}(\theta) = \hat{\mathbb{E}}_t \left[\min\!\Big(
        \rho_t(\theta)\,\hat{A}_t,\;
        \text{clip}\big(\rho_t(\theta),\, 1{-}\varepsilon,\, 1{+}\varepsilon\big)\,\hat{A}_t
    \Big)\right]
\end{equation}
where $\rho_t(\theta) = \frac{\pi_\theta(a_t \mid s_t)}{\pi_{\theta_{\text{old}}}(a_t \mid s_t)}$ is the probability ratio between the current and reference policy.

\paragraph{Key insight: the reference policy is not a separate network.} During each rollout collection phase, the policy generates actions and stores the corresponding log-probabilities $\log \pi_{\theta_{\text{old}}}(a_t \mid s_t)$. These scalar values are the ``reference''. During the subsequent $K$ gradient epochs over the same batch, the current network recomputes $\log \pi_\theta(a_t \mid s_t)$, and the ratio $\rho_t$ measures how much the policy has drifted. At epoch 1, $\rho_t = 1$ exactly; by epoch $K$, it may have drifted, and clipping prevents it from exceeding $[1{-}\varepsilon,\, 1{+}\varepsilon]$.

The advantage $\hat{A}_t$ is computed via GAE:
\begin{equation}
    \hat{A}_t = \sum_{l=0}^{H-t-1} (\gamma \lambda_{\text{GAE}})^l \,\delta_{t+l}, \qquad
    \delta_t = r_t + \gamma V_\phi(s_{t+1}) - V_\phi(s_t),
\end{equation}
where $V_\phi$ is the value function (critic) sharing a backbone with the policy. GAE with $\lambda_{\text{GAE}} = 0.95$ provides a good bias--variance trade-off: it looks multiple steps ahead without accumulating too much noise from the value function estimate.

The full loss combines three terms:
\begin{equation}
    L = L^{\text{CLIP}} - c_v \, L^{\text{VF}} + c_e \, H[\pi_\theta]
\end{equation}
where $L^{\text{VF}} = (V_\phi(s_t) - V_t^{\text{target}})^2$ is the value function loss, $H[\pi_\theta]$ is the policy entropy, and $c_v, c_e$ are weighting coefficients.

\subsection{Training Cycle}

Each PPO update follows this cycle:
\begin{enumerate}[nosep]
    \item \textbf{Collect}: run the current policy for $n_{\text{steps}} = 50$ steps across $n_{\text{envs}} = 14$ parallel environments, producing $50 \times 14 = 700$ transitions per update.
    \item \textbf{Compute advantages}: apply GAE with $\gamma=0.99$, $\lambda_{\text{GAE}}=0.95$.
    \item \textbf{Optimise}: for $K = 10$ epochs, shuffle the 700 transitions into mini-batches of 256, compute the clipped surrogate loss, and update via Adam ($\alpha = 3 \times 10^{-4}$).
    \item \textbf{Discard buffer} and return to step 1 with the updated policy.
\end{enumerate}

The 14 parallel environments (one per CPU core) provide two benefits: (i) each gradient update uses $14\times$ more data, reducing variance; (ii) wall-clock time is the same as a single environment since the rollouts are parallelised via \texttt{SubprocVecEnv}.

\subsection{Hyperparameter Choices and Justification}

\begin{table}[H]
\centering
\begin{tabular}{llp{7.5cm}}
\toprule
Parameter & Value & Rationale \\
\midrule
\texttt{learning\_rate} & $3 \times 10^{-4}$ & Standard PPO default; works well for MLPs \\
\texttt{n\_steps} & 50 & Equals episode length; each rollout is one full episode per env \\
\texttt{batch\_size} & 256 & 256/700 $\approx$ 37\% of buffer per mini-batch; good bias--variance balance \\
\texttt{n\_epochs} & 10 & 10 passes allow thorough optimisation of each batch; clipping prevents overfitting \\
\texttt{gamma} & 0.99 & High discount: reward at step 49 is worth $0.99^{49} \approx 0.61$ of step 0 \\
\texttt{gae\_lambda} & 0.95 & Moderate lookahead in advantage estimation \\
\texttt{clip\_range} & $0.2 \to 0.0$ & Linear schedule: allows large policy changes early (exploration), forces small updates late (exploitation) \\
\texttt{ent\_coef} & 0.0 & See discussion below \\
\texttt{log\_std\_init} & $-1.0$ & Initial std $= e^{-1} \approx 0.37$; default of 1.0 causes excessive exploration \\
\texttt{net\_arch} & $[128, 128]$ & Two hidden layers with ReLU; sufficient for $N=4$ \\
\texttt{VecNormalize} & reward only & Running reward normalisation with clip at 10.0 \\
\bottomrule
\end{tabular}
\caption{PPO hyperparameters and their justification.}
\label{tab:hyperparams}
\end{table}

\subsubsection{Why We Dropped the Entropy Bonus ($c_e = 0$)}

In standard PPO, a positive entropy coefficient $c_e > 0$ encourages exploration by penalising low-entropy (deterministic) policies. In our setting, \textbf{the entropy bonus was actively harmful} for three reasons:

\begin{enumerate}[nosep]
    \item \textbf{Unbounded action std inflation}: with $c_e > 0$, the entropy term $H[\pi] = \frac{1}{2}\log(2\pi e \sigma^2)$ is maximised by increasing $\sigma$. Since our Gaussian policy has a learnable $\log \sigma$, the entropy bonus pushes $\sigma$ upward without bound. Actions become random noise rather than structured shipping decisions.

    \item \textbf{Reward hacking}: a high-std policy samples near-zero actions on average (since $\mathbb{E}[\max(0, a)] \to 0.5\sigma$ as $\sigma \to \infty$ under the action clipping). The environment then interprets this as ``ship nothing'', yielding zero transport cost. The policy converges to a local optimum where it avoids all shipping and accepts the demand penalty.

    \item \textbf{Exploration is not the bottleneck}: with 14 parallel environments seeing different demand realisations, and with initial $\sigma = 0.37$ providing adequate coverage, the policy has enough diversity to explore the action space without an explicit entropy incentive.
\end{enumerate}

Instead of entropy, we control the exploration--exploitation trade-off through:
\begin{itemize}[nosep]
    \item \textbf{Clip range schedule} ($0.2 \to 0.0$): allows large policy updates early (exploration), then shrinks to near-zero (exploitation). This naturally transitions the policy from broad search to fine-tuning.
    \item \textbf{Low initial log-std} ($-1.0$): the policy starts with $\sigma \approx 0.37$, which produces meaningful actions from the beginning rather than random noise.
\end{itemize}

\subsubsection{Exploration vs.\ Exploitation in Detail}

The exploration--exploitation dynamics in this problem are subtle. The action space is $\mathbb{R}^{N^2} = \mathbb{R}^{16}$ (for $N=4$), and the effective action after projection and clipping is constrained to the feasible polytope $\{T \geq 0,\, T_{ii}=0,\, \sum_j T_{ij} \leq I_i\}$. Random exploration in this high-dimensional space is extremely inefficient: most random actions correspond to either shipping nothing (negative raw values, clipped to 0) or shipping everything (which depletes inventory in one step).

A well-initialised policy with $\sigma=0.37$ concentrates most probability mass on small positive actions (ship a little), which is a much better starting point for gradient-based optimisation. As training progresses and the clip range shrinks, the policy commits to the discovered structure without reverting to random exploration.

%% ======================================================================
\section{Implementation Details}
%% ======================================================================

\subsection{Architecture}

The codebase consists of three main modules:

\begin{enumerate}[nosep]
    \item \texttt{env/warehouse\_env.py}: core MDP logic — state transitions, reward computation, action projection, and demand sampling.
    \item \texttt{env/warehouse\_gym\_env.py}: thin Gymnasium wrapper that normalises observations and maps the $[-1,1]^{N^2}$ action space to transport quantities via $T = \text{clip}(a, 0, 1) \cdot I_{\max}$.
    \item \texttt{train\_sb3.py}: training script using SB3's PPO with parallel environments, reward normalisation, and WandB logging.
\end{enumerate}

\subsection{Action Scaling}

A critical design choice is how the network's raw output $a \in [-1, 1]^{N^2}$ maps to shipping quantities. Our final mapping is:
\[
    T_{ij} = \text{clip}(a_{ij},\, 0,\, 1) \cdot I_{\max}
\]
This ensures that $a=0$ (the network's natural initialisation) maps to \textbf{zero shipping}, which is a safe default. An earlier version used $(a+1)/2 \cdot I_{\max}$, where $a=0$ mapped to $I_{\max}/2 = 250$ units --- causing the policy to start by shipping massive quantities, then learn to ``do nothing'' as a local optimum (reward hacking).

\subsection{Reward Normalisation}

We apply \texttt{VecNormalize} from SB3 to maintain a running estimate of reward variance and normalise each reward by dividing by $\hat{\sigma}_r$. Combined with \texttt{clip\_reward=10}, this prevents outlier rewards from destabilising the value function. Observation normalisation is unnecessary because the obs vector is already hand-normalised to $[0, 1]$.

\subsection{Computational Considerations}

Training runs on CPU rather than GPU. For the small MLP architecture ($[128, 128]$, $\sim$35k parameters), GPU kernel launch overhead dominates computation time. With 14 parallel environments on a 14-core CPU (AMD Ryzen, 4.2 GHz), 1M timesteps complete in approximately 15 minutes.

%% ======================================================================
\section{Results}
%% ======================================================================

\subsection{PPO Learning Curves}

Over 2M timesteps (approximately 1,400 evaluation episodes on 14 parallel environments):

\begin{table}[H]
\centering
\begin{tabular}{lcccc}
\toprule
Metric & Start (ep 100) & End (ep 1400) & Trend \\
\midrule
Transport cost & $\sim$800 & $\sim$10 & $\downarrow$ converged \\
Episode reward & $-$16,000 & $-$7,000 & $\uparrow$ improving \\
Demand satisfaction & 55\% & 78\% & $\uparrow$ plateau \\
Inventory std & 400 & 180 & $\downarrow$ improved \\
\bottomrule
\end{tabular}
\caption{PPO training progress on the small scenario ($N=4$, seed 42).}
\end{table}

\paragraph{Suggested figures for the poster:}
\begin{itemize}[nosep]
    \item \textbf{Figure 1}: Four-panel learning curves (transport cost, episode reward, demand satisfaction, inventory std) vs.\ evaluation step. This shows the convergence behaviour and plateau.
    \item \textbf{Figure 2}: Shipping graph at ep 550 showing which warehouse pairs are connected. This reveals whether the policy discovered the hub structure.
    \item \textbf{Figure 3}: Greedy baseline comparison --- overlay greedy (flat line) with PPO curves to show that PPO matches or approaches heuristic performance.
\end{itemize}

\subsection{Comparison with Greedy Baseline}

\begin{table}[H]
\centering
\begin{tabular}{lccc}
\toprule
Agent & Demand Satisfaction & Transport Cost & Episode Reward \\
\midrule
Random & $\sim$30\% & High & $\sim -60{,}000$ \\
Greedy & $\sim$74\% & $\sim$6{,}000 & $\sim -42{,}000$ \\
PPO (1M steps) & $\sim$78\% & $\sim$10 & $\sim -7{,}000$ \\
\bottomrule
\end{tabular}
\caption{Agent comparison on small scenario (seed 42, with reward normalisation).}
\end{table}

PPO achieves \textbf{comparable demand satisfaction to the greedy heuristic} while discovering a policy with dramatically lower transport cost. The greedy baseline ships aggressively (high transport cost, high satisfaction), while PPO learns a more efficient redistribution strategy.

%% ======================================================================
\section{Why Convergence Is Difficult}
%% ======================================================================

The demand satisfaction plateau at $\sim$75--78\% is not a failure of the algorithm but rather reflects fundamental properties of the problem:

\begin{enumerate}
    \item \textbf{Stochastic demand floor}: with Poisson demand ($\mu=10$) and finite inventory, there is an irreducible probability that demand exceeds available stock at some warehouse regardless of shipping decisions. The theoretical upper bound on satisfaction with optimal redistribution is well below 100\%.

    \item \textbf{High-dimensional continuous action space}: the action is $\mathbb{R}^{N^2}$ (16 dimensions for $N=4$), projected onto a polytope. Unlike discrete problems where $Q$-learning can enumerate actions, PPO must learn a continuous mapping from a 9-dimensional observation to a 16-dimensional action. The ratio of action dimensions to observation dimensions ($16/9 \approx 1.8$) is unusually high for RL problems, making credit assignment difficult.

    \item \textbf{Delayed and diffuse reward signal}: a shipping decision at step $t$ affects inventory at step $t{+}1$, which affects whether demand is met at step $t{+}1$. With $\gamma = 0.99$ and $H=50$, the temporal credit assignment chain is long. Moreover, the reward sums over all $N$ warehouses, so the signal for any single $T_{ij}$ is diluted.

    \item \textbf{Asymmetric cost structure}: the hub-and-spoke cost matrix creates a combinatorial structure where routing through hubs is cheap but direct spoke-to-spoke transfer is expensive. Discovering this routing structure purely from scalar reward feedback (without any graph-based inductive bias) is non-trivial.

    \item \textbf{Greedy is near-optimal}: the greedy baseline, which simply ships surplus to deficit, already achieves $\sim$74\% satisfaction. This suggests that the myopic (single-step) solution is close to the multi-step optimum, leaving little room for RL to improve through temporal planning. The remaining gap to 100\% is dominated by demand stochasticity, not by suboptimal planning.
\end{enumerate}

These observations suggest that a \textbf{graph neural network (GNN) policy} with explicit graph structure (hub-and-spoke topology encoded as edge features) could break through the plateau by providing the inductive bias that a flat MLP lacks.

%% ======================================================================
\section{Conclusion}
%% ======================================================================

\begin{enumerate}[nosep]
    \item PPO with Stable Baselines 3 learns a shipping policy that achieves 78\% demand satisfaction and near-zero transport cost, matching the greedy heuristic on service level while being significantly more cost-efficient.

    \item Reward design is critical: normalising each reward component to $[0,1]$ before weighting prevents any single term from dominating the gradient signal. Without normalisation, the policy collapses to ``ship nothing'' (reward hacking).

    \item The entropy bonus, standard in PPO, is counterproductive in this setting because it inflates the action standard deviation in a continuous-action problem with projection constraints. Replacing it with a clip-range schedule and low initial std provides more controlled exploration.

    \item The 75--78\% demand satisfaction plateau is likely an inherent limit of (a) stochastic demand with finite inventory and (b) flat MLP policies without structural knowledge of the hub-and-spoke topology.

    \item \textbf{Extensions}: a GNN policy that encodes the cost matrix as graph edge features, combined with a $\lambda$-sensitivity sweep to characterise the cost--service trade-off frontier, would be the most impactful additions.
\end{enumerate}

\begin{thebibliography}{9}
\bibitem{schulman2017ppo}
J.~Schulman, F.~Wolski, P.~Dhariwal, A.~Radford, and O.~Klimov,
``Proximal Policy Optimization Algorithms,'' \textit{arXiv:1707.06347}, 2017.
\end{thebibliography}

\end{document}
